\section{Related work}
\label{sec:related}

\subsection{Knowledge management and document acquisition}
Studies of LLM integration with knowledge-based methods describe both the
potential and the challenges of combining language models with structured
knowledge \citep{yang-kbs-llm-2025}. Public knowledge bases organise facts under explicit ontologies
\citep{auer-dbpedia-kb-2007,hoffart-yago2-kb-2013}, whereas enterprise content
management emphasises governance over organisational documents
\citep{paivarinta-ecm-2005,tyrvaainen-ecm-2006}. Research on organisational
knowledge explains how tacit experience becomes shareable
\citep{nonaka-seci-1994,nonaka-seci-ba-2000}; adjacent work studies metadata
quality, tagging, human review and auditability
\citep{park-metadata-quality-2009,golder-tagging-2006,mosqueira-rey-hitl-2022,barredo-arrieta-xai-2020}; automated text categorization provides the classification backbone \citep{sebastiani2002machine}. These strands motivate reliable placement
and governance, but do not by themselves connect those controls to evidence-
checked answers. Document-AI methods, including layout-aware analysis and
MinerU, recover text and structure from complex documents
\citep{zhong-publaynet-layout-2019,huang-layoutlmv3-docai-2022,mineru}. Near-
duplicate detection and provenance provide complementary safeguards during
ingestion \citep{manku-near-dup-2007,missier-prov-2013,cheney-provenance-2009}.
HMRS treats acquisition quality, automatic classification and traceable
indexing as design commitments to be tested rather than as preprocessing
assumptions; the campaigns measure where they hold and where they fail,
without establishing a causal link to answer quality.

A line of work in this journal takes an explicitly architectural form: a
software reference architecture for journalistic knowledge platforms is
derived from the literature and a prototype, specifying module
responsibilities and architecture principles \citep{ocaa2023software}, and an
earlier workshop version outlines the architecture and a proof-of-concept
\citep{gallofr2021jkpdesign}. Other work builds and evaluates complete
knowledge-grounded systems---dialogue systems with graph-based semantic
models \citep{yang2024building}, OpenQA systems for low-resource languages
\citep{budur2024building}, an LLM-based clinical decision support system
\citep{vito2025heliot}, and an IR system for bank supervisors
\citep{aarab2024system}---while privacy-preserving service frameworks
\citep{liao2026practical} and LLM-assisted construction of expert systems
\citep{garridomer2026gofai} treat deployment constraints as design inputs.
Domain knowledge bases are built either from operator--function descriptions
\citep{hou2025method} or by instantiating ontologies with LLM-generated domain
knowledge \citep{ciatto2025large}. The same systems orientation appears in an
open-source multi-agent framework for evidence-based knowledge-graph enrichment
\citep{mohamed2026orchestra}, LLM-agent systems for data-driven engineering
modelling \citep{liu2026automatin}, end-to-end LLM systems for circuit analysis
\citep{chen2026enhancing}, and an intelligent system that automates
hyperparameter tuning for gradient-boosted trees \citep{sorokin2023sigopt}.
HMRS shares this systems orientation but
makes collection membership, derived descriptions and lifecycle state the
auditable objects, rather than proposing a reference architecture or a domain
ontology.

\subsection{Retrieval, organisation and agentic search}
RAG grounds generated responses in retrieved passages \citep{rag-lewis}; dense
retrievers learn shared query--passage representations \citep{dpr}, while BM25
remains a strong lexical baseline \citep{bm25}. Ghali et al. combine
semantic retrieval, vector storage and language models to query structured and
unstructured data \citep{ghali-kbs-retrieval-2025}. Hybrid retrieval combines
sparse and dense signals, including reciprocal-rank fusion
\citep{rrf-cormack}; embedding models such as BGE-M3 support multilingual and
multi-granularity search \citep{bge-m3}. Graph- and tree-structured systems
organise and retrieve structured relationships: GraphRAG adds entity-relation
graphs and community summaries \citep{graphrag}, LightRAG uses dual-level
retrieval with incremental graph merging \citep{lightrag}, and RAPTOR builds a
recursive cluster tree \citep{raptor}. RMR-SST represents multimodal
engineering-document content as structured semantic trees and retrieves
across abstraction levels \citep{li-kbs-rmrsst-2026}. Hierarchical
representation is therefore already an established retrieval concern: RAPTOR summarises a tree of clusters over
document text and retrieves over that tree, and GraphRAG retrieves over an
entity graph induced from the corpus. HMRS differs in what the hierarchy
governs. Its hierarchy is a management structure over collections and
documents---which collection exists, which document belongs to it, which
description explains it---and selected lifecycle outcomes can be audited
across its state sources. The hierarchies that RAPTOR and GraphRAG query are built
offline at index time---recursive cluster summarisation and pre-generated
community summaries respectively---and their papers do not address
incremental maintenance after corpus change \cite{raptor,graphrag}. LightRAG
does contribute incremental graph updates \citep{lightrag}, so HMRS's
difference is not incremental update per se but the governance of collection
membership and document lifecycle, reconciled across three state sources.
Incremental vector-index maintenance is itself an
active systems concern \cite{singh2021freshdiskann,xu2023spfresh}; HMRS
complements it at the metadata layer with an explicit re-index operation and
three-source reconciliation. Retrieval granularity is likewise a design
variable: proposition-level and late-chunking strategies change what a
retrievable unit carries \cite{chen2024dense,gunther2024late}. TEIE-RAG
combines hierarchical semantic chunking with retrieval-model tuning and source
grounding for engineering information extraction \citep{shen-kbs-teierag-2026}.
HMRS splits at recognised structural markers and indexes the resulting parts.
HMRS uses a
narrower document--library--tag graph for navigation and candidate
expansion; the graph does not establish answer evidence on its own.

The decision of which collection to consult has a longer history in
distributed IR: resource-selection methods such as CORI and ReDDE rank
collections by estimated relevance before searching within them
\citep{callan1995searching,si2003redde}. HMRS adopts the same
select-then-search ordering but treats collections as managed, reconciled
state rather than static statistical summaries.

Tool-using agents extend retrieval through iterative interaction
\citep{react-yao,toolformer}. MCP standardises access to tool servers
\citep{mcp}, while Modular RAG and agentic-retrieval surveys describe
reconfigurable retrieval pipelines and their efficiency challenges
\citep{modular-rag,singh-agentic-rag}. DeepRead performs structure-aware
reasoning for agentic search: it converts a document to structured Markdown,
indexes paragraphs with section coordinates, and lets an agent locate then
read sections in order \citep{deepread}. Its hierarchy is thus a read-only,
per-document navigation aid built once at ingestion. HMRS differs in scope
and in lifecycle: its hierarchy spans collections and documents, is a
persistent managed asset whose state can be audited after selected operations, and exposes the
catalogue-first librarian lane alongside dense and hybrid alternatives under
explicit call and tool policies.

Recent articles in this venue extend the retrieval stack in complementary
directions: retrieval-augmented generation over engineering design knowledge
\citep{siddharth2024retrieval}, soft-prompt tuning of dense retrievers
\citep{peng2024soft}, and an evaluation of semantic-relation query expansion
in retrieval systems \citep{massai2023evaluatio}. Private-library retrieval
supports code generation \citep{zan2025private}; work on knowledge graphs
ranges from studying how LLMs comprehend graph facts supplied in different
prompt formats \citep{dai2025large} to prompting LLMs with graphs for
long-tail factual questions \citep{huang2025prompting}, while fine-grained
stateful exploration navigates graphs through an LLM \citep{tao2025fine}.
Real-time financial data grounding \citep{sinha2026finbloom} and coordinated
multi-agent answer generation \citep{saadaoui2025coordinat} illustrate the
range of deployed configurations. These systems optimise what is retrieved
and how it is used; HMRS instead asks which managed collection is
authoritative and whether its state can be audited after each lifecycle
operation.

\subsection{Verification, correction and abstention}
Two further lines border the design without competing with it. Agent memory
systems treat memory as an agent-managed store of interaction records:
MemGPT pages memory between a context window and an external store
\citep{memgpt}, Mem0 extracts and consolidates long-term memory from
conversations \citep{mem0}, and A-MEM organises notes as an agentic memory
graph \citep{amem}; continual-learning work adds an information-theoretic dual
memory system \citep{wu2025informati}. HMRS manages a document corpus with audited lifecycle
state; the stores are complementary rather than competing, and an ingest
agent of the kind these frameworks assume can also use collection assignment
and consistency checks of the kind examined here. Knowledge organisation
standards such as SKOS \citep{skos} and Dublin Core provide conventions for
vocabularies and descriptive metadata; HMRS derives collection and document
descriptions from parsed content at ingest time.

Self-RAG and CRAG integrate evidence assessment or retrieval correction into
generation \citep{selfrag,crag}; FLARE triggers retrieval during generation
\citep{flare}. Selective prediction under domain shift \citep{kamath},
identifying what LLMs do not know \citep{feng-abstain,yin-knowwhat},
ontology-grounded retrieval \citep{cyberbot}, and citation verifiability
\citep{liu-verifiability,gao-citations} address adjacent concerns. EPRAG
diagnoses evidence states to choose responses, clarification or abstention in
multi-source enterprise retrieval \citep{kinalioglu-kbs-eprag-2026}. HMRS externalises
verification as a reusable server-side operation: the gate fetches document
text, judges evidence and returns provenance to the answer layer. This design
makes the verifier inspectable, but not infallible.

The LLM-as-a-judge literature documents preference bias, imperfect agreement,
listwise instability and self-preference \citep{zheng-llmjudge,sun-urara,panickssery-selfpref}; outcome-based verification can reward flawed reasoning
\citep{lightman-verify}. These findings motivate separating evidence
verification from rank ordering and testing abstention at the answer layer.

Trustworthiness and evaluation are also studied at the systems level in this
venue: a systematic meta-survey of explainable AI \citep{saeed2023explainab},
logic-grounded metamorphic testing of reasoning reliability
\citep{zhou2026lgmt}, and a benchmark plus evaluation framework for multi-step
clinical diagnosis \citep{hou2025msdiagnos}. These contributions supply the
vocabulary of reliability and graded evaluation; the distinction claimed here
is narrower: HMRS reports the measured failure mode of its own evidence gate
and turns that negative result into a configuration rule.

\subsection{Research gap}
\label{sec:gap}

Supplementary Table S11 compares representative strands rather than every
system in the literature. Prior work already studies selective question
answering, citation verifiability and retrieval correction; the distinction
here is not that abstention or evidence checks are absent from prior research.
The gap this study addresses is therefore a systems question: how can a
knowledge-document platform make the management state that precedes retrieval
observable, maintainable and usable by an agent? HMRS examines the complete
workflow from acquisition checks and automatic collection assignment through
three-source consistency, lifecycle operations, catalogue-first access and
answer-time evidence inspection. It reports where implementation and transport
boundaries constrain the workflow and measures verifier admission separately
from verifier ordering. The contribution is an auditable management--retrieval
contract, not a claim that prior systems lack reusable verification or that one
retrieval algorithm is universally superior.

The study treats three design commitments as questions to test: whether
content-based classification and lifecycle checks expose collection state;
whether a retrieved document should be re-ranked by a verifier or retain its
fusion order; and whether agent workflows can report absence while separating
completed negative searches from execution failures. Its contribution is a
knowledge document management workflow that connects and exposes stages from
parse admission through answer-evidence reporting, a per-question account of
answer-layer checks, and a same-pool diagnostic showing that ordering by the
verifier score, not threshold filtering, is where the measurable loss sits. It
proposes no new ranking algorithm and does not claim that the underlying
retrieval problem is solved.
